\documentclass[11pt]{article}
\usepackage[margin=1in]{geometry}
\usepackage[T1]{fontenc}
\usepackage[utf8]{inputenc}
\usepackage{amsmath,amssymb,amsthm}
\usepackage{enumitem}

\title{ICPC World Finals 2025\\I. Slot Machine}
\author{}
\date{}

\begin{document}
\maketitle

\section*{Problem Summary}

The machine has $n$ wheels, all carrying the same cyclic order of $n$ symbols. Rotating a wheel only
changes its offset in that common cycle. After every action we are told only the number of distinct
symbols currently visible. We must use at most $10\,000$ actions to make all wheels show the same
symbol.

\section*{Key Observations}

\begin{itemize}[leftmargin=*]
    \item Fix all wheels except one wheel $i$, and rotate only $i$ through a full cycle. For each position of
    wheel $i$, the reported number of distinct symbols is minimal exactly when the symbol shown by
    wheel $i$ is already present on at least one of the other wheels. So one full scan of a wheel gives a
    binary mask of which symbols are currently present elsewhere.
    \item First scan wheel $1$. From this mask we can pick one symbol $s$ that is absent from all other
    wheels, and one symbol $r$ that is present on some other wheel.
    \item Now fix wheel $1$ to show $s$ and scan another wheel $j$. Then fix wheel $1$ to show $r$ and
    scan $j$ again. The only difference between the two scans is the position where wheel $j$ itself
    shows symbol $s$. Therefore comparing the two masks reveals exactly how much wheel $j$ must be
    rotated to show $s$.
    \item Repeating this for every wheel gives a common target symbol, so the final phase is just rotating
    every wheel to its discovered offset.
\end{itemize}

\section*{Algorithm}

\begin{enumerate}[leftmargin=*]
    \item Maintain the current offset of each wheel modulo $n$.
    \item Define \texttt{scan(i)}:
    rotate wheel $i$ through all $n$ positions, record the reported distinct-count value for each one, and
    mark the positions that attain the minimum value.
    \item Scan wheel $1$ once. Choose one position labeled $0$ in the resulting mask and call its symbol
    $s$; choose one position labeled $1$ and call its symbol $r$.
    \item Rotate wheel $1$ to $s$. For each wheel $j \ge 2$:
    \begin{itemize}[leftmargin=*]
        \item scan $j$ with wheel $1$ fixed at $s$;
        \item rotate wheel $1$ to $r$ and scan $j$ again;
        \item compare the two masks to find the unique position where wheel $j$ shows symbol $s$;
        \item rotate wheel $1$ back to $s$.
    \end{itemize}
    \item Finally, rotate every wheel to the stored position that makes it show $s$.
\end{enumerate}

\section*{Correctness Proof}

We prove that the algorithm returns the correct answer.

\paragraph{Lemma 1.}
For a scan of wheel $i$, a position is marked by the algorithm if and only if the symbol shown by wheel
$i$ at that position is already shown by at least one other wheel.

\paragraph{Proof.}
Fix the other $n-1$ wheels. If wheel $i$ is rotated to a symbol that is absent from all other wheels, the
total number of distinct visible symbols increases by one compared to showing any already-present
symbol. Therefore the minimum possible distinct-count value is attained exactly at the positions where
wheel $i$ shows a symbol already present elsewhere. \qed

\paragraph{Lemma 2.}
For every wheel $j \ge 2$, comparing the two scans described above identifies the unique offset at which
wheel $j$ shows the chosen target symbol $s$.

\paragraph{Proof.}
The two scans differ only in the symbol shown by wheel $1$: once it is $s$, once it is $r$.
All other wheels are unchanged. By construction, $s$ is absent from the other wheels when wheel $1$
shows $s$, while $r$ is present on at least one other wheel when wheel $1$ shows $r$.
Hence the membership test from Lemma 1 changes exactly at the position where wheel $j$ itself shows
$s$. That position is therefore identified uniquely. \qed

\paragraph{Theorem.}
The algorithm outputs the correct answer for every valid input.

\paragraph{Proof.}
By Lemma 2, for every wheel we discover an offset that makes it show the same symbol $s$.
After the final rotation phase, all wheels therefore display $s$, so the machine shows only one distinct
symbol and the jackpot condition is met. \qed

\section*{Complexity Analysis}

One scan uses exactly $n$ actions. We perform one scan for wheel $1$, two scans for every other wheel,
and then $O(n)$ final rotations. Thus the number of actions is at most
\[
1 + n + 2(n-1)n + O(n) < 2n^2 + 3n \le 5150
\]
for $n \le 50$, well below the limit of $10\,000$. The local memory usage is $O(n)$.

\section*{Implementation Notes}

\begin{itemize}[leftmargin=*]
    \item The helper \texttt{rotateTo} always uses the shorter cyclic rotation.
    \item The solution exits immediately when the judge reports $k=1$, as required by the interaction
    protocol.
\end{itemize}

\end{document}
